\documentclass[../main.tex]{subfiles}
\begin{document}
\chapter{2013-2014学年微积分（一）（上）期中考试}

\section{基本计算题(每小题 6 分，共 60 分)}
\begin{enumerate}
    \item 计算数列极限$l=\lim_{n\to\infty}\frac{1}{4}\cdot\frac{3}{8}\cdot\frac{5}{12}
              \cdots\frac{2n-1}{4n}$.
          \vspace{11em}

    \item 计算极限$l=\lim_{x\to 0}\frac{2^{x^2}-3^{x^2}}{\left(2^{x}-3^{x}\right)^{2}}$.
          \vspace{11em}

    \item 计算极限$l=\lim_{x\to 0}\frac{\mathrm{e}^{\tan x}-\mathrm{e}^{\sin x}}{x^{3}}$.
          \vspace{11em}

    \item 设$f(x)=(\sin x)^{\cos x}$，求$f'(x)$.
          \vspace{10em}

    \item 设参数方程$\begin{cases}
                  x=t+t^{2}, \\
                  y=t\mathrm{e}^{t}
              \end{cases}$确定了函数$y=y(x)$，求$\frac{\mathrm{d}y}{\mathrm{d}x}$，$\frac{\mathrm{d}^{2}
                  y}{\mathrm{d}x^{2}}$.
          \vspace{10em}

    \item 设函数$y=y(x)$由方程$y=h(x^{2}+y^{2})$所确定，$h$处处可导，且$h'\neq\frac{1}{2y}$，求$y
              '$.
          \vspace{10em}

    \item 求函数$f(x)=\sin\ln(1+x)$在$x=0$处带Peano余项的3阶Taylor公式.
          \vspace{10em}

    \item 求函数$f(x)=\frac{x}{1+2x}$在$x=0$处的n阶导数$f^{(n)}(0)$.
          \vspace{10em}

    \item 求曲线$y=x\ln\left(\mathrm{e}+\frac{1}{x}\right)$的斜渐近线方程.
          \vspace{13em}

    \item 指出函数$f(x)=\left(1-\mathrm{e}^{\frac{x}{1-x}}\right)^{-1}$的间断点与类型.
          \vspace{13em}
\end{enumerate}

\section{综合题(每小题 7 分，共 28 分)}
\begin{enumerate}
    \setcounter{enumi}{10}
    \item 设函数$f(x)$在$x=0$处连续，$f(0)\neq0$，且对一切$x,y$有$f(x+y)=f(x
              )\cdot f(y)$，证明$f(x)$处处连续.
          \vspace{13em}

    \item 讨论方程$\ln x-\frac{x}{\mathrm{e}}+1=0$有几个实根？并给出证明.
          \vspace{12em}

    \item 设函数$f(x)$在$x=2$处可导，且$f(2)\neq0$.求$\lim_{x\to\infty}\left
              [\frac{f\left(2+\frac{1}{x}\right)}{f(2)}\right]^{x}$.
          \vspace{12em}

    \item 有一个顶点朝上的圆锥形容器，高为$8\mathrm{m}$，底半径为$R=2\sqrt{2}
              \mathrm{m}$，向其中注水.设当水深$h=6\mathrm{m}$时，水面上升的速度为$\frac{\mathrm{d}h}{\mathrm{d}t}
              =\frac{4}{\pi}\mathrm{m/min}$，求此时水的体积的变化率$\frac{\mathrm{d}V}{\mathrm{d}t}$.
          \vspace{11em}
\end{enumerate}

\section{证明题(每小题 6 分，共 12 分)}
\begin{enumerate}
    \setcounter{enumi}{14}
    \item 证明不等式：当$0<x<1$时，$\frac{1-x}{1+x}<\mathrm{e}^{-2x}$.
          \vspace{12em}

    \item 设$0<x_{1}<x_{2}$，证明$x_{1}\ln x_{2}-x_{2}\ln x_{1}=(\ln\xi-1)(x
                  _{1}-x_{2})$，其中$\xi$介于$x_{1}$与$x_{2}$之间.
          \vspace{12em}
\end{enumerate}

\chapter{2013-2014学年微积分（一）（上）期中考试参考答案}
\section{基本计算题(每小题 6 分，共 60 分)}
\begin{enumerate}
    \item \textbf{Solution}. 设$x_{n}=\frac{1}{4}\cdot\frac{3}{8}\cdot\frac{5}{12}
              \cdots\frac{2n-1}{4n}$，则
          \[
              0<x_{n}=\frac{2n-1}{4n}x_{n-1}<\frac{1}{2}x_{n-1}<\frac{1}{2^{2}}x_{n-2}<
              \cdots<\frac{1}{2^{n-1}}x_{1},
          \]
          且$\lim_{n\to\infty}\frac{1}{2^{n-1}}x_{1}=0$，由夹逼定理得$l=\lim_{n\to\infty}
              x_{n}=0$.

    \item \textbf{Solution}.
          \[
              \begin{aligned}
                  l = {} & \lim_{x\to 0}\frac{\mathrm{e}^{x^2\ln 2}-\mathrm{e}^{x^2\ln 3}}{\left(\mathrm{e}^{x\ln 2}-\mathrm{e}^{x\ln 3}\right)^2}                                                   \\
                  = {}   & \lim_{x\to0}\frac{1+x^2\ln 2-1-x^2\ln 3+o(x^2)}{\left(1+x\ln 2-1-x\ln 3+o(x)\right)^2}=\lim_{x\to0}\frac{(\ln 2 - \ln 3)x^2+o(x^2)}{\left[(\ln 2 - \ln 3)x+o(x)\right]^2} \\
                  = {}   & \lim_{x\to0}\frac{(\ln 2-\ln 3)x^2+o(x^2)}{(\ln 2-\ln 3)^2x^2+o(x^2)}                                                                                                     \\
                  = {}   & \frac{1}{\ln 2-\ln 3}.                                                                                                                                                    \\
              \end{aligned}
          \]

    \item \textbf{Solution}. 由Lagrange中值定理，存在$\xi$介于$\tan x$与$\sin x$之间，使得$\mathrm{e}
              ^{\xi}(\tan x-\sin x)=\mathrm{e}^{\tan x}-\mathrm{e}^{\sin x}$.

          当$x\to0$时，$\tan x\to0$，$\sin x\to0$，故$\xi\to0$.所以
          \[
              \begin{aligned}
                  l= {} & \lim_{x\to0}\mathrm{e}^{\xi}\cdot\frac{\tan x-\sin x}{x^3}                                                      \\
                  = {}  & \lim_{x\to0}\frac{\sec^2 x-\cos x}{3x^2}=\lim_{x\to0}\frac{\sec^2 x- 1}{3x^2}+\lim_{x\to0}\frac{1-\cos x}{3x^2} \\
                  = {}  & \frac{1}{3}+\frac{1}{6}=\frac{1}{2}.
              \end{aligned}
          \]

    \item \textbf{Solution}. 因为$\ln f(x) =\cos x\ln\sin x$，所以
          \[
              f'(x)=\left(\sin x
              \right)^{\cos x}\left(-\sin x\ln\sin x+\cos x\cdot\frac{\cos x}{\sin x}\right
              ).
          \]

    \item \textbf{Solution}.
          \[
              \frac{\mathrm{d}y}{\mathrm{d}x}=\frac{\frac{\mathrm{d}y}{\mathrm{d}t}}{\frac{\mathrm{d}x}{\mathrm{d}t}}
              =\frac{(t+1)\mathrm{e}^{t}}{2t+1},
          \]
          \[
              \frac{\mathrm{d}^{2}y}{\mathrm{d}x^{2}}=\frac{\mathrm{d}}{\mathrm{d}t}\left
              (\frac{\mathrm{d}y}{\mathrm{d}x}\right)\cdot\frac{\mathrm{d}t}{\mathrm{d}x}
              =\left(\frac{(t+1)\mathrm{e}^{t}}{2t+1}\right)'\cdot\frac{1}{2t+1}=\frac{(t+2)\mathrm{e}^{t}\cdot(2t+1)-2(t+1)\mathrm{e}^{t}}{(2t+1)^{3}}
              =\frac{(2t+3)t\mathrm{e}^{t}}{(2t+1)^{3}}.
          \]

    \item \textbf{Solution}. 方程$y=h(x^{2}+y^{2})$两边对$x$求导得
          \[
              y'=h'(x^{2}+y^{2})(2x+2yy'),
          \]
          整理得$y'=\frac{2xh'(x^{2}+y^{2})}{1-2yh'(x^{2}+y^{2})}$.

    \item \textbf{Solution}. $\ln(1+x)=x-\frac{x^{2}}{2}+\frac{x^{3}}{3}+o(x^{3})$，所以
          \[
              \begin{aligned}
                  f(x) = {} & \sin\ln(1+x)=\sin\left(x-\frac{x^2}{2}+\frac{x^3}{3}+o(x^{3})\right)                                                                                                                         \\
                  = {}      & \left(x-\frac{x^2}{2}+\frac{x^3}{3}+o(x^{3})\right)-\frac{1}{6}\left(x-\frac{x^2}{2}+\frac{x^3}{3}+o(x^{3})\right)^{3}+o\left(\left(x-\frac{x^2}{2}+\frac{x^3}{3}+o(x^{3})\right)^{3}\right) \\
                  = {}      & x-\frac{x^2}{2}+\frac{x^3}{3}+o(x^{3})-\frac{x^3}{6}+o(x^{3})                                                                                                                                \\
                  = {}      & x-\frac{x^2}{2}+\frac{x^3}{6}+o(x^{3}).
              \end{aligned}
          \]

    \item \textbf{Solution}. 因为$f(x)=\frac{1}{2}-\frac{1}{2}\cdot\frac{1}{2x+1}$，所以$f
                  ^{(n)}(x)=-\frac{1}{2}\cdot\frac{(-1)^{n}\cdot n!\,2^{n}}{(2x+1)^{n+1}}$，

          因此
          \[
              f^{(n)}(0)=(-1)^{n-1}\cdot n!\cdot 2^{n-1}.
          \]

    \item \textbf{Solution}. 因为
          \[
              \lim_{x\to\infty}\frac{x\ln\left(\mathrm{e}+\frac{1}{x}\right)}{x}
              =1,\quad\lim_{x\to\infty}x\ln\left(\mathrm{e}+\frac{1}{x}\right)-x=\lim_{t\to0}
              \frac{\ln(\mathrm{e}+t)-1}{t}=\lim_{t\to0}\frac{1}{\mathrm{e}+t}= \frac{1}{\mathrm{e}},
          \]
          所以曲线$y=x\ln\left(\mathrm{e}+\frac{1}{x}\right)$的斜渐近线方程为$y=x+\frac{1}{\mathrm{e}}$.

    \item \textbf{Solution}. 函数$f(x)$的间断点为$x=0$和$x=1$.

          当$x\to0$时，$\mathrm{e}^{\frac{x}{1-x}}\to1$，所以$\lim_{x\to0}f(x)=\lim_{x\to0}
              \left(1-\mathrm{e}^{\frac{x}{1-x}}\right)^{-1}=\infty$，故$x=0$为无穷间断点.

          当$x\to1^{+}$时，$\mathrm{e}^{\frac{x}{1-x}}\to0$，所以$\lim_{x\to1^+}f(x)=
              \lim_{x\to1^+}\left(1-\mathrm{e}^{\frac{x}{1-x}}\right)^{-1}=1$；

          当$x\to1^{-}$时，$\mathrm{e}^{\frac{x}{1-x}}\to\infty$，所以$\lim_{x\to1^-}
              f(x)=\lim_{x\to1^-}\left(1-\mathrm{e}^{\frac{x}{1-x}}\right)^{-1}=0$，故$x=
              1$为跳跃间断点.
\end{enumerate}

\section{综合题(每小题 7 分，共 28 分)}
\begin{enumerate}
    \setcounter{enumi}{10}
    \item \textbf{Solution}. 在方程$f(x+y)=f(x)\cdot f(y)$中令$x=y=0$得$f(0)
              =f^{2}(0)$，因$f(0)\neq0$，所以$f(0)=1$.

          任取$x_{0}\in\mathbf{R}$，在方程$f(x+y)=f(x)\cdot f(y)$中令$x=x_{0}$，得$f(
              x_{0}+y)=f(x_{0})\cdot f(y)$.令$y\to0$，得
          \[
              \lim_{y\to0}f(x_{0}+y)=f(x_{0})\cdot\lim_{y\to0}f(y)=f(x_{0})\cdot f(0)=f
              (x_{0}),
          \]
          所以$f(x)$在$x=x_{0}$处连续.由$x_{0}$的任意性可知$f(x)$处处连续.

    \item \textbf{Solution}. 设$f(x)=\ln x-\frac{x}{\mathrm{e}}+1$，则$f'(x)
              =\frac{1}{x}-\frac{1}{\mathrm{e}}$，$f(x)$在$(0,\mathrm{e})$上单调递增，在$(
              \mathrm{e},+\infty)$上单调递减.

          因为$f(\mathrm{e})=1>0$，$f(0^{+})=-\infty$，$f(\mathrm{e}^{3})=4-\mathrm{e}
              ^{2}<0$，所以方程$\ln x-\frac{x}{\mathrm{e}}+1=0$有两个实根.

    \item \textbf{Solution}.
          \[
              \begin{aligned}
                  \lim_{x\to\infty}\left[\frac{f\left(2+\frac{1}{x}\right)}{f(2)}\right]^{x} = {} & \exp\left\{{\lim\limits_{x\to\infty}x\left[\ln f\left(2+\frac{1}{x}\right)-\ln f(2)\right]}\right\} \\
                  = {}                                                                            & \exp\left[{\lim\limits_{t\to0}\frac{\ln f(2+t)-\ln f(2)}{t}}\right]                                 \\
                  = {}                                                                            & \exp\left[{\left(\ln f(t)\right)'\Big|_{t=2}}\right]=\exp\left({\frac{f'(2)}{f(2)}}\right).
              \end{aligned}
          \]

    \item \textbf{Solution}. 由几何关系，$V=\frac{1}{3}\cdot 8 \cdot 8\pi-\frac{1}{3}
              \pi\cdot\frac{(8-h)^{3}}{8}$.方程两边对$t$求导得
          \[
              \frac{\mathrm{d}V}{\mathrm{d}t}=\frac{1}{3}\pi\cdot\frac{3(8-h)^{2}}{8}\cdot
              \frac{\mathrm{d}h}{\mathrm{d}t}=\frac{\pi(8-h)^{2}}{8}\cdot\frac{\mathrm{d}h}{\mathrm{d}t}
              .
          \]
          将$h=6\mathrm{m}$，$\frac{\mathrm{d}h}{\mathrm{d}t}=\frac{4}{\pi}\mathrm{m/min}$代入上式得$\frac{\mathrm{d}V}{\mathrm{d}t}
              =2\mathrm{m^3/min}$.
\end{enumerate}

\section{证明题(每小题 6 分，共 12 分)}
\begin{enumerate}
    \setcounter{enumi}{14}
    \item \textbf{Proof}. 取对数，即证$\ln\frac{1-x}{1+x}<-2x(0<x<1)$.

          令$f(x)=\ln\frac{1-x}{1+x}+2x$，则$f'(x)=-\frac{2}{1-x^{2}}+2=-\frac{2x^{2}}{1-x^{2}}
              <0$，所以$f(x)$在$(0,1)$上单调递减.

          因为$f(0)=0$，所以当$0<x<1$时，$f(x)<f(0)=0$，即$\frac{1-x}{1+x}<\mathrm{e}
              ^{-2x}$.

    \item \textbf{Proof}. 设$f(x)=\frac{\ln x}{x}$，$g(x)=\frac{1}{x}$，则函数$f
              (x),g(x)$在$[x_{1},x_{2}]$上连续，在$(x_{1},x_{2})$内可导，且$g(x)\neq0$.

          由Cauchy中值定理，存在$\xi$介于$x_{1}$与$x_{2}$之间，使得
          \[
              \frac{f(x_{2})-f(x_{1})}{g(x_{2})-g(x_{1})}=\frac{\frac{\ln x_{2}}{x_{2}}-\frac{\ln
                      x_{1}}{x_{1}}}{\frac{1}{x_{2}}-\frac{1}{x_{1}}}=\frac{\frac{\frac{1}{\xi}\cdot
                      \xi-\ln\xi}{\xi^{2}}}{-\frac{1}{\xi^{2}}}=\ln\xi-1,
          \]
          化简得$x_{1}\ln x_{2}-x_{2}\ln x_{1}=(\ln\xi-1)(x_{1}-x_{2})$.
\end{enumerate}
\end{document}



